\section{USAMO 1981}
        \begin{enumerate}
        	\item (\textit{USAMO 1981}) Prove that if $n$ is not a multiple of $3$, then the angle $\frac{\pi}{n}$ can be trisected with ruler and compasses. 


 



\textbf{Solution}

Let $n=3k+1$.  Multiply throughout by $\pi/3n$. We get 

 $\frac{\pi}{3} = \frac{\pi \times k}{n} + \frac{\pi}{3n}$

 Re-arranging, we get

 $\frac{\pi}{3} - \frac{\pi \times k}{n} = \frac{\pi}{3n}$

 A way to interpret it is that if we know the value $k$, then the remainder angle of subtracting $k$ times the given angle from $\frac{\pi}{3}$ gives us $\frac{\pi}{3n}$, the desired trisected angle.

 This can be extended to the case when $n=3k+2$ where now, the equation becomes$\frac{\pi}{3} - \frac{\pi \times k}{n} = \frac{2\pi}{3n}$

 Hence in this case, we will have to subtract $k$ times the original angle from $\frac{\pi}{3}$ to get twice the the trisected angle. We can bisect it after that to get the trisected angle.

 \textbf{Generalization} If regular polygons of $m$ sides and $n$ sides can be constructed, where $m$ and $n$ are relatively prime integers greater than or equal to three, then regular polygons of $mn$ sides can be constructed. Indeed, such a polygon can be constructed by first constructing an $m$-gon, and then creating $m$ distinct $n$-gons with at least one vertex being a vertex of the $m$-gon.

 


	\item (\textit{USAMO 1981}) What is the largest number of towns that can meet the following criteria. Each pair is directly linked by just one of air, bus or train. At least one pair is linked by air, at least one pair by bus and at least one pair by train. No town has an air link, a bus link and a train link. No three towns, $A, B, C$ are such that the links between $AB, AC$ and $BC$ are all air, all bus or all train.


 



\textbf{Solution}

Assume $AB$, $AC$, and $AD$ are all rail.

 None of $BC$, $CD$, or $CD$ can be rail, as those cities would form a rail triangle with $A$.

 If $BC$ is bus, then $BD$ is bus as well, as otherwise $B$ has all three types.

 However, $CD$ cannot be rail (as $\triangle ACD$ would be a rail triangle), bus (as $BCD$ would be a bus triangle), or ferry (as $C$ and $D$ would have all three types).

 Therefore, no city can have three connections of the same type.

 
Assume there are 5 towns - $A$, $B$, $C$, $D$, and $E$.

 Two connections from $A$ must be of one type, and two of another; otherwise there would be at least three connections of the same type from $A$, which has been shown to be impossible.

 Let $AB$ and $AC$ be rail connections, and $AD$ and $AE$ be bus.

 Assume $CD$ is air.

 $BC$ cannot be rail ($\triangle ABC$ would be a rail triangle) or bus ($C$ would have all three types), so $BC$ must be air.

 $DE$ cannot be bus ($\triangle ADE$ would be a bus triangle) or rail ($D$ would have all three types), so $DE$ must be air.

 $BE$ cannot be rail ($E$ would have all three types) or bus ($B$ would have all three types), so $BE$ must be air.

 However, $BD$ cannot be rail ($D$ would have all three types), bus ($B$ would have all three types), or air ($D$ would have three air connections).

 Therefore, the assumption that $CD$ is air is false.

 $CD$ can equally be rail or bus; assume it is bus.

 $BC$ cannot be rail ($\triangle ABC$ would be a rail triangle) or air ($C$ would have all three types), so $BC$ must be bus.

 $BD$ cannot be air ($B$ would have all three types) or bus ($D$ would have three bus connections), so $BD$ must be rail.

 $DE$ cannot be air ($D$ would have all three types) or bus ($D$ would have three bus connections), so $DE$ must be rail.

 $CE$ cannot be air ($C$ would have all three types) or bus ($E$ would have three bus connections), so $CE$ must be rail.

 The only connection remaining is $BE$, which cannot be orange as both $B$ and $D$ would have all three types, but this means there are no air connections.

 Therefore, it is impossible with five (or more) towns.

 
A four-town mapping is possible:

 $AB$, $BC$, $CD$, and $DA$ are connected by bus.

 $AC$ is connected by rail.

 $BD$ is connected by air.

 Therefore, the maximum number of towns is $4$.

 


	\item (\textit{USAMO 1981}) Show that for any triangle, $\frac{3\sqrt{3}}{2}\ge \sin(3A) + \sin(3B) + \sin (3C) \ge -2$. 


 When does the equality hold?


 



\textbf{Solution}

Given three angles that add to $180^\circ$, one can construct a triangle from them. However, its angles must all be nonnegative; thus the constraints on the angles of a triangle are $0^\circ\le A,B,C\le180^\circ$ and $A+B+C=180^\circ$.

 In fact, at this point, we only care about $3A, 3B,$ and $3C$.  Let us call them $x, y,$ and $z$.  We have:

 
 \[0 \le x,y,z \le 540^\circ, x+y+z = 540^\circ\]

 and we must prove that

 
 \[\frac{3\sqrt{3}}{2}\ge \sin(x) + \sin(y) + \sin (z) \ge -2\]

 Without loss of generality, assume $x \le y \le z$.  It follows that $x \le 180^\circ$ and $z \ge 180^\circ$ (otherwise, $x+y+z$ would be strictly greater than $180^\circ$ or strictly less than $180^\circ$, respectively).

 Since $\sin(u)$ is nonnegative over the interval $[0^\circ, 180^\circ]$ and $-1 \le \sin (u) \le 1$ for all real $u$, we can immediately use the $x \le 180^\circ$ result to show that

 
 \[\sin( x) + \sin (y) + \sin (z) \le 0 + -1 + -1 = -2\]

 This proves the lower bound; equality occurs when $\sin(x) = 0$ and $\sin (y) = \sin (z) = -1$, and this is reachable only when $y = z = 270^\circ$ and $x = 0^\circ$, which translates into $A = 0^\circ$ and $B = C = 90^\circ$.

 Now for the upper bound. It is true that $\sin(u)$ is non-positive over the interval $[180^\circ, 360^\circ]$.  Also, $\frac{3\sqrt{3}}{2}>2$ (by a simple squaring argument). Then $z \ge 180^\circ$.  If $z \le 360^\circ$, then $\sin (z) \le 0$, and then
 \[\sin (x) + \sin (y) + \sin (z) \le 1 + 1 + 0 = 2 < \frac{3\sqrt{3}}{2}\]Therefore, we need to handle the case where $z \ge 360^\circ$. $\sin(u + 360^\circ) = \sin(u)$, so we can subtract $360^\circ$ from $z$ due to periodicity. Then the constraints become $0^\circ \le x,y,z \le 180^\circ$ and $x+y+z = 180^\circ$.

 Let us suppose that $(x,y,z)$ is the selection of $x,y,z$ given the above constraints with the largest value of $\sin (x) + \sin (y) + \sin (z)$. We shall see that, given these constraints, this value increases when the distance between two of these variables decreases, and it is maximized when $x=y=z$.

 $\textbf{Lemma:}$ If $f(x)$ is well-behaved and $f'(x)$ is strictly decreasing over an open interval $(a, b)$, then, for any $x,y$ selected from that interval with $x \ne y$, then $2f\left(\frac{x + y}{2}\right) > f(x) + f(y).$

 $\textbf{Proof of lemma:}$ It is true that, for any well-behaved function $f(x)$, $f(x+d) = f(x) + \int_{x}^{x+d} f'(u)\,du$, and likewise, $f(x-d) = f(x) - \int_{x-d}^x f'(u)\,du$. Without loss of generality, suppose $x < y$; let $d = \frac{1}{2}(y-x)$.  Since $f'(u)$ is given to be strictly decreasing over the relevant interval, we have $f'(u) > f'(x+d)$ for all $u\in(x, x+d)$; therefore
 \[f(x+d) = f(x) + \int_x^{x+d} f'(u)\,du > f(x) + \int_x^{x+d} f'(x+d)\,du = f(x) + d\cdot f'(x+d)\]Likewise, $f'(u) < f'(y-d)$ for all $u\in(y-d, y)$; therefore
 \[f(y-d) = f(y) - \int_{y-d}^y f'(u)\,du > f(y) - \int_{y-d}^y f'(y-d)\,du = f(y) - d\cdot f'(y-d)\]Therefore, adding our inequalities together,
 \[f(x+d) + f(y-d) > f(x) + d\cdot f'(x+d) + f(y) - d\cdot f'(y-d)\]But $x+d = y-d = \frac{x+y}{2}$. In that case, the $d\cdot f'(x+d)$ and $-d\cdot f'(y-d)$ terms cancel, and we get our desired result: $2 f\left(\frac{x+y}{2}\right) > f(x) + f(y)$.

 Since $\sin(u)$ is well-behaved, and $\sin'(u) = \cos(u)$, which is strictly decreasing over the open interval $(0^\circ, 180^\circ)$, then we can apply our lemma. Suppose for the sake of contradiction that $x,y,z$ are not all the same value, and this produces a maximum. Then, without loss of generality, suppose $x \ne y$. Let $a = \frac{x+y}{2}$. Then the selection $(a,a,z)$ also satisfies the constraints $0^\circ \le a,a,z \le 180^\circ$ and $a+a+z \le 180^\circ$. Furthermore, by the above lemma, $2\cdot\sin (a) > \sin (x) + \sin (y)$, so this new selection $(a,a,z)$ has a larger corresponding value $\sin (a) + \sin (a) + \sin (z)$ than the selection $(x,y,z)$. This contradicts our original assumption that $(x,y,z)$ was chosen to have the maximum. Therefore, $x=y=z$.

 Then, since $x+y+z = 180^\circ$, we conclude $x=y=z=60^\circ$. Then
 \[\sin (x) + \sin (y) + \sin (z) = 3 \cdot \sin (60^\circ) = 3 \cdot \frac{\sqrt{3}}{2}\]So $\frac{\sqrt{3}}{2}$ is indeed the maximum value of $\sin(x) + \sin (y) + \sin (z)$, proving the bound. Translating back into the previous problem, the maximum occurs when $x=y=60^\circ$ and $z=420^\circ$; translating into the original problem, we have proven the upper bound, and equality with the upper bound occurs when $A=B=20^\circ$ and $C=140^\circ$.

 


	\item (\textit{USAMO 1981}) The sum of the measures of all the face angles of a given complex polyhedral angle is equal to the sum of all its dihedral angles. Prove that the polyhedral angle is a trihedral angle.


 



\textbf{Solution}

$n=3$ is certainly possible. For example, take $\angle APB = \angle APC = \angle BPC = 90^\circ$ (so that $PA$, $PB$, $PC$ are mutually perpendicular). Then the three planes through P are also mutually perpendicular, so the two sums are both $270^\circ$.

 We show that $n>3$ is not possible.

 It is almost obvious that the sum of the $n$ dihedral angles at P is less than $360^\circ$. Take another plane which meets the lines $PA$, $PB$, $PC$ etc at $A'$, $B'$, $C'$, ... and so that the foot of the perpendicular from $P$ to the plane lies inside the $n$-gon $A'B'C'\cdots$ then as we move $P$ down the perpendicular the angles $A'PB'$ etc all increase. But when it reaches the plane their sum is $360^\circ$. Take any point O inside the $n$-gon $ABC\cdots$ . We have $\angle PBA + \angle PBC > \angle ABC$. Adding the n such equations we get $\sum(180^\circ - APB) > \sum ABC = (n - 2) 180^\circ$. So, $\sum APB... < 360^\circ$The sum of the $n$ angles between the planes is at least $180(n-2)^\circ$. If we take a sphere center $P$, then the lines $PA$, $PB$ intersect it at $A", B", ...$ which form a spherical polygon. The angles of this polygon are the angles between the planes. We can divide the polygon into $n - 2$ triangles. The angles in a spherical triangle sum to at least $180^\circ$. So the angles in the spherical polygon are at least $180(n-2)^\circ$. So we have $(n - 2) 180^\circ < 360^\circ$ and hence $n < 4$.

 


	\item (\textit{USAMO 1981}) Show that for any positive real $x$, 
 \[[nx]\ge \sum_{k=1}^{n}\left(\frac{[kx]}{k}\right),\]
where $[x]$ denotes the greatest integer not exceeding $x$.


 



\textbf{Solution(due to Pavel Zatitskiy)}

First of all we write $[kx]=kx-\{kx\}$. So, we need to prove that$\sum_{1}^{n}\left(\frac{\{kx\}}{k}\right)\geq \{nx\}.$Let's denote $a_k=\{kx\}$. It is easy to see that $a_k+a_m \geq a_{k+m}$. We need to prove$\sum_{1}^{n}\left(\frac{a_k}{k}\right)\geq a_n.$

 We will prove it by induction by $n$. The base is obvious, so we need to make a step.

 Let's take $m$ such that $\frac{a_m}{m}$ is minimal. If $m=n$ then our inequality is obvious. So, $m<n$. Then, by induction, $\sum_{1}^{n-m}\left(\frac{a_k}{k}\right)\geq a_{n-m}$ and $\sum_{n-m+1}^{n}\left(\frac{a_k}{k}\right)\geq m\cdot \frac{a_m}{m}=a_m$. Now we can add these two inequalities and get
 \[\sum_{1}^{n}\left(\frac{a_k}{k}\right) \geq a_{n-m}+a_m \geq a_n\]

 


\end{enumerate}